\documentclass[a4paper, 12pt]{exam}

\usepackage[utf8]{inputenc}
\usepackage[T1]{fontenc}
\usepackage[french]{babel}
\usepackage{lmodern}
\usepackage{amsmath}
\usepackage{geometry}
\usepackage{xcolor}
\geometry{margin=2.5cm}

% Personnalisation de l'en-tête et du pied de page
\pagestyle{headandfoot}
\firstpageheader{IFFD}{Travaux Dirigés - Cadastre}{Contentieux des taxes foncière}
\runningheader{Travaux Dirigés}{}{Contentieux des taxes foncière}
\firstpagefooter{}{Page \thepage\ sur \numpages}{}
\runningfooter{}{Page \thepage\ sur \numpages}{}

% Personnalisation de l'affichage des solutions
\renewcommand{\solutiontitle}{\noindent\textbf{Éléments de correction :}\par\noindent}
\unframedsolutions % ou \framedsolutions pour encadrer les réponses
\SolutionEmphasis{\itshape\color{blue}} % Nécessite le package xcolor si utilisé, sinon juste \itshape

\begin{document}
	
	\begin{center}
		\Large \textbf{CAS PRATIQUES : INSTRUCTION DES RÉCLAMATIONS} \\[0.5cm]
	\end{center}
	
	\vspace{0.5cm}
	
\section*{ANALYSE DE CONTENTIEUX}

Analyser les cas ci-dessous au regard : de la forme des réclamations, des délais de réclamation, et des décisions à prendre. Pour chacun des cas proposés vous indiquerez la solution à retenir en motivant vos décisions (articles du CGI...) et la suite à donner pour clôturer l'affaire (admission totale, admission partielle, rejet).

\begin{questions}
	
	\question Un propriétaire qui a acquis sa maison depuis un an environ, vous envoie un courrier indiquant qu'il estime son imposition trop élevée. Il demande le mode de calcul de sa taxe.
	\begin{solution}
		Le propriétaire ne conteste pas son imposition, il demande seulement des détails sur le mode de calcul. La demande sera traitée en « réponse directe ». 
		NB : à l'issue de la réponse, il arrive que le propriétaire dépose une réclamation.
	\end{solution}
	
	\question Une maison d'habitation, commencée en 2019, a été achevée le 20/08/2020. Elle est édifiée sur une parcelle de 11 ares, classée depuis 2018 en terrain à bâtir. Le propriétaire sollicite un dégrèvement de la taxe foncière 2020 sur les propriétés non bâties. Il joint à sa demande l'avis d'imposition de 2020 et la déclaration concernant sa maison. La demande parvient au service le 30/12/2020 et est instruite le 16/02/2021.
	\begin{solution}
		Contestation de la TFPNB pour 2020 (évaluation).
		\begin{itemize}
			\item \textbf{En la forme} : Réclamation recevable pour l'année 2020 (dans les délais jusqu'au 31/12/2021).
			\item \textbf{Au fond} : il s'agit d'une construction neuve, imposable en bâti à compter du 01/01/2021. Imposition en TAB pour 2020 justifiée (construction non achevée au 01/01/2020). Article 270 du CGI : principe de l'annualité.
			\item \textbf{DECISION} : REJET. Rédaction et notification du rejet au contribuable.
			\item Déclaration hors délai : perte de l'exonération pour 2021 mais exonération pour 2022. Article 266 du CGI.
		\end{itemize}
	\end{solution}
	
	\question Le 18/12/2020 M LAWSON adresse une réclamation au service portant sur la taxe foncière de l'année 2020. Il conteste l'imposition de sa boulangerie, suite à sa cessation d'activité en octobre 2019. Il précise que le local est inutilisé depuis cette date. Les services des impôts confirment la cessation d'activité au 15/10/2019. L'instruction du dossier est effectuée le 8/01/2021.
	\begin{solution}
		Contestation de la TFB pour 2020 (évaluation).
		\begin{itemize}
			\item \textbf{En la forme} : Réclamation recevable pour l'année 2020 (dans les délais jusqu'au 31/12/2021).
			\item \textbf{Au fond} : contestation de l'imposition au motif que la cessation d'activité est intervenue le 15/10/2019. Article 270 du CGI Annualité de l'impôt : le local existe au 01/01/2020. Article 87 du LPF : Le local n'a pas changé d'affectation, même s'il y a cessation d'activité.
			\item \textbf{DECISION} : REJET. Rédaction et notification du rejet au contribuable.
			\item NB : Maintien de l'évaluation du local pour la TF.
		\end{itemize}
	\end{solution}
	
	\question Le 02/11/2020 la STE NORFLEX adresse une réclamation de taxe foncière. Elle s'étonne de ne pas bénéficier de l'exonération de 2 ans  pour l'année 2020, alors que son bâtiment est neuf. Le service qui instruit le dossier constate le dépôt d'une déclaration avec une date d'achèvement de travaux donnée par le contribuable au 01/02/2019. La déclaration a été rédigée le 18/06/2019 et reçue au service le 21/06/2019. L'instruction du dossier est effectuée le 15/12/2020.
	\begin{solution}
		Contestation de la TFB pour 2020 (exonération).
		\begin{itemize}
			\item \textbf{En la forme} : Réclamation recevable au titre de l'année 2020 (dans les délais jusqu'au 31/12/2021).
			\item \textbf{Au fond} : Contestation sur l'application de l'exonération de 2 ans  pour une construction neuve (art 262 du CGI). Article 266 du CGI : l'exonération est subordonnée à la souscription de la déclaration dans les 04 mois de la D.A.T. Achèvement des travaux : 01/02/2019. Dépôt de la déclaration : 21/06/2019. Article 266 : déclaration hors délais. L'exonération s'applique pour la période restant à courir après le 31/12 de l'année suivante.
			\item \textbf{DECISION} : REJET. Rédaction et notification du rejet au contribuable.
			\item NB : Imposition année 2020, exonération année 2021.
		\end{itemize}
	\end{solution}
	
	%		\question Un verger a été arraché en octobre 2019. Le terrain a été ensemencé aussitôt en luzerne. Le propriétaire demande un dégrèvement pour 2020, le tarif des vergers étant supérieur à celui des terres. Le service n'a pas eu connaissance de cette nouvelle affectation. La demande est reçue le 12/11/2020, accompagnée de l'avis d'imposition de l'année 2020. L'instruction du dossier est effectuée le 20/01/2021.
	%		\begin{solution}
		%			Contestation de la TFPNB pour 2020 (évaluation).
		%			\begin{itemize}
			%				\item \textbf{En la forme} : Réclamation recevable au titre de l'année 2020 (dans les délais jusqu'au 31/12/2021).
			%				\item \textbf{Au fond} : changement d'affectation soumis au système déclaratif ; défaut de déclaration, mais constatation annuelle des changements par le service. Question de fait avec avis au maire. Articles 1406 et 1517 du CGI et R* 198-3 du LPF.
			%				\item \textbf{DECISION} : ADMISSION TOTALE. Mise à jour de la base par le calcul non-bâti (4B-3D). Dégrèvement prononcé si supérieur ou égal à 8 € (Article 1965 L du CGI).
			%			\end{itemize}
		%		\end{solution}
	
	\question Par lettre en date du 15 octobre 2020, Mme Durand vous indique qu'il lui est impossible de régler sa taxe foncière en raison du décès de son mari et des frais d'obsèques qu'elle a payés.
	\begin{solution}
		La personne ne demande pas l'application d'une règle de droit tendant à diminuer son impôt : elle ne conteste pas mais demande une mesure de tempérance. Il s'agit d'une demande de remise gracieuse qui sera examinée par le service compétent pour son impôt sur le revenu.
	\end{solution}
	
	\question Le 6 janvier 2021, M. LANGEVIN a adressé une demande, portant sur l'année 2020. Il conteste l'évaluation assignée à son hôtel particulier. Le rôle de l'année 2020 a été mis en recouvrement le 31/08/2020. L'agent terrain constate que le classement en 3ème catégorie des locaux d'habitation est fondé et qu'une erreur a été commise dans le calcul de la surface pondérée du local. L'instruction, assurée le 10 février 2021, permet d'établir que la rectification de l'erreur de surface ramène la valeur locative de 5 870 200 FCFA à 5 396 000 FCFA.
	\begin{solution}
		Contestation de la TFB 2020 (évaluation).
		\begin{itemize}
			\item \textbf{En la forme} : Réclamation recevable au titre de l'année 2020 (dans les délais jusqu'au 31/12/2021).
			\item \textbf{Au fond} : le classement en catégorie 3 est maintenu ; le calcul de la surface pondérée est revu, ce qui implique une diminution de VL. Article 270 du CGI.
			\item \textbf{DECISION} : ADMISSION PARTIELLE. Motif du rejet : classement.
			\item Dégrèvement 2020 et mise à jour pour 2021 de la base par le calcul bâti. Mise en œuvre du pouvoir de restitution (Art. 418 du LPF).
		\end{itemize}
	\end{solution}
	
	\question Le 10/01/2021, Maître RENOIR, avocat, saisit le service, pour contester les impositions à la taxe foncière des années 2020 et 2019 de la société QUICAMPOIX. Celui-ci demande que soit appliquée pour ces années l'exonération de 2 ans pour construction neuve. La consistance du local précisée sur la déclaration donne : Terrain 156 233 000 F et Bâtiment 1 587 225 F. L'exonération a été effectuée sur le bâtiment uniquement.
	\begin{solution}
		Contestation de la TFB pour les années 2019 et 2020 (exonération).
		\begin{itemize}
			\item \textbf{En la forme} : La réclamation reçue le 10/01/2021 est recevable uniquement au titre de l'année 2020 (délais jusqu'au 31/12/2021) ; hors délai au titre de l'année 2019 (délais jusqu'au 31/12/2020).
			\item \textbf{Au fond} : Contestation sur l'application de l'exonération de 2 ans pour une construction neuve (art 262 du CGI). Article 266 du CGI : l'exonération est subordonnée à la souscription de la déclaration dans les 04 mois de la D.A.T. La déclaration a été souscrite dans les délais, l'exonération doit s'appliquer au local (bâtiment et terrain).
			\item \textbf{DECISION} : ADMISSION TOTALE POUR 2020. Mise en œuvre du pouvoir de restitution (article 418 du LPF) : Degrièvement d'office pour 2019.
		\end{itemize}
	\end{solution}
	
	\question M. ARNAL a produit le 12/11/2020 une déclaration pour sa maison d'habitation. Le service a considéré cette maison comme terminée en avril 2019, une imposition supplémentaire a été établie pour 2020. M. ARNAL conteste et démontre que sa construction a été achevée le 18/08/2020. L'enquête confirme ses dires.
	\begin{solution}
		Contestation de la TFB 2020 (évaluation).
		\begin{itemize}
			\item \textbf{En la forme} : Réclamation recevable pour 2020.
			\item \textbf{Au fond} : la construction achevée le 18/08/2020, est imposable à compter du 01/01/2021. La déclaration est dans les délais. Exonération de droit commun de 5 ans applicable. Pour 2020, imposition en terrain à bâtir. Articles 270 du CGI.
			\item \textbf{DECISION} : ADMISSION TOTALE. Dégrèvement de l'imposition supplémentaire.
		\end{itemize}
	\end{solution}
	
	\question M DEMURGER adresse le 09/09/2020 une réclamation concernant sa taxe foncière 2020. Il conteste la nouvelle valeur locative servant de base au calcul suite à l'intégration de la révision (contestation du tarif, demande de classement en secteur 3 au lieu de 4).
	\begin{solution}
		Contestation TFB pour l'année 2020 (évaluation).
		\begin{itemize}
			\item \textbf{En la forme} : La réclamation est recevable (dans les délais jusqu'au 31/12/2021).
			\item \textbf{Au fond} : Le secteur et le tarif sont des paramètres de la révision. Contestation impossible.
			\item \textbf{DECISION} : REJET.
		\end{itemize}
	\end{solution}
	
\end{questions}



\section*{ANALYSE DE CAS  - ATTRIBUTION}

\begin{questions}
	\question Le 17 décembre 2020, M DURAND adresse une réclamation par laquelle il indique qu'il n'est plus propriétaire d'une parcelle. Il a été imposé pour celle-ci au titre de l'année N-1 alors que la vente a été réalisée le 15/09/2019. L'acte a été publié au livre foncier le :
	\begin{itemize}
		\item 4 octobre 2019,
		\item 20 janvier 2020,
		\item la publication n'est pas encore intervenue
	\end{itemize}
	\begin{solution}
		Contentieux d'attribution. Contestation de la TF 2020.
		\begin{itemize}
			\item \textbf{Cas a) et b)} : Réclamation recevable en la forme. Mise en œuvre de l'article 284 du CGI. Dégrèvement du propriétaire imposé à tort et imposition du redevable légal. \textbf{DECISION : ADMISSION TOTALE}.
			\item \textbf{Cas c)} : Réclamation recevable en la forme. \textbf{DECISION : REJET} (article 283 du CGI).
		\end{itemize}
	\end{solution}
	
	\question Par acte en date du 10/10/2019, publié au livre foncier le 17/11/2019, M AUGERE a vendu la parcelle à M MAURY. Un échec de la liaison automatique Livre foncier/Matrice Cadastrale  est édité le 26/06/2020.
	\begin{solution}
		Erreur d'attribution pour 2020, décelée après l'envoie de la matrice pour le recouvrement de l'année 2020 : affaire d'office. 
		\textbf{DECISION : DEGREVEMENT D'OFFICE POUR 2020 POUR LE VENDEUR ET IS (Imposition Supplémentaire ) POUR 2020 POUR L'ACQUEREUR}.
	\end{solution}
	
	\question Le 02/06/2020, M DURAND adresse une réclamation indiquant qu'il a vendu la parcelle le 18/08/2018 et demande le remboursement de la TF 2019 et la mise à jour pour 2020. L'acte a été publié le 10/10/2018.
	\begin{solution}
		Contentieux d'attribution. Contestation de la TF 2019. Réclamation recevable. Article 1404 du CGI.
		\begin{itemize}
			\item \textbf{Pour l'année 2019} : Dégrèvement du propriétaire imposé à tort et imposition du redevable légal. \textbf{DECISION : ADMISSION TOTALE}.
			\item \textbf{Pour l'année 2020} : Mise à jour de la matrice et imposition de l'acquéreur. 
		\end{itemize}
	\end{solution}
\end{questions}


\section*{ANALYSE DE CAS  - IMPOSITIONS SUPPLÉMENTAIRES}

\begin{questions}
	\question Le 15 décembre 2020, en consultant le plan au service, M. ROSEDONC a fait observer que son garage de 60 M2, construit pourtant en 2011 n'était pas encore représenté. Il n'a été ni déclaré, ni évalué.
	\begin{solution}
		Changement de consistance (addition de construction). Défaut de déclaration pour une propriété bâtie (Article 262 du CGI).
		\begin{itemize}
			\item La sous-évaluation doit être réparée par la mise en œuvre d'un rôle particulier (RP).
			\item Le changement est de 2011 : on retiendra pour le RP la base et les taux 2020 avec coefficient multiplicateur 5 (Art. 266 CGI), avec mise à jour pour l'année 2021.
		\end{itemize}
	\end{solution}
	
	\question M. ROLLINA a construit une piscine, achevée le 12 mars 2017 (renseignements obtenus auprès des services de l'urbanisme). Il n'a jamais adressé de déclaration au service.\\
	
	Suite à une exploitation de photographies aériennes, l'agent terrain a constaté, le 12 novembre 2020 :
	\begin{itemize}
		\item la présence effective d'une piscine en bois semi-enterrée,
		\item  la transformation du garage en pièce d'habitation depuis le 15 mars 2019
	\end{itemize}
	\begin{solution}
		\begin{itemize}
			\item \textbf{1er changement (piscine)} : Changement de consistance, défaut du système déclaratif réparé par un rôle particulier (Art 266 du CGI ). RP base et taux 2020x3, mise à jour pour 2021.
			\item \textbf{2ème changement (garage)} : Changement de caractéristiques physiques, réparé par un rôle supplémentaire pour 2020, avec mise à jour pour 2021.
		\end{itemize}
	\end{solution}
	
\end{questions}

	
\end{document}